
%%%%%%%%%%
%% 2019 %%
%%%%%%%%%%

@article{Malani2019trialSubjects,
   author = {Malani, Anup and Philipson, Tomas J.},
   title = {Labor Markets in Statistics: The Subject Supply Effect in Medical R&D},
   journal = {Journal of Human Capital},
   volume = {13},
   number = {2},
   pages = {293-340},
   note = {doi: 10.1086/702923},
   abstract = {Medical R&D differs from other R&D because of a unique linkage between output and input markets: potential consumers of existing medical products are also potential subjects in clinical trials required to develop new products. Therefore, a quality increase or price reduction for an existing treatment reduces patients? incentive to participate in trials of new treatments. We label this linkage the ?subject supply effect? and provide evidence of it from trials of hepatitis C and AIDS treatments. The subject supply effect has important positive implications for how policies affect the rate of medical R&D and normative implications for whether subjects ought to be compensated for enrolling in clinical trials.},
   ISSN = {1932-8575},
   DOI = {10.1086/702923},
   url = {https://doi.org/10.1086/702923},
   year = {2019},
   type = {Journal Article}
}

%%%%%%%%%%
%% 2017 %%
%%%%%%%%%%

@article{BaudeMalani2017doctrinalMethods,
   author = {Baude, William and Chilton, Adam S. and Malani, Anup},
   title = {Making Doctrinal Work More Rigorous: Lessons from Systematic Reviews Symposium: Developing Best Practices for Legal Analysis},
   journal = {University of Chicago Law Review},
   volume = {84},
   number = {1},
   pages = {37-58},
   url = {https://heinonline.org/HOL/P?h=hein.journals/uclr84&i=43},
   year = {2017},
   type = {Journal Article}
}

@article{BaudeMalani2017postiveLegalMethodology,
   author = {Baude, William and Chilton, Adam S. and Malani, Anup},
   title = {A Call for Developing a Field of Positive Legal Methodology Symposium: Developing Best Practices for Legal Analysis},
   journal = {University of Chicago Law Review},
   volume = {84},
   number = {1},
   pages = {1-6},
   url = {https://heinonline.org/HOL/P?h=hein.journals/uclr84&i=7},
   year = {2017},
   type = {Journal Article}
}


@article{LakdawallaMalani2017insuranceValueInnovation,
   author = {Lakdawalla, Darius and Malani, Anup and Reif, Julian},
   title = {The insurance value of medical innovation},
   journal = {Journal of Public Economics},
   volume = {145},
   pages = {94-102},
   abstract = {Economists think of medical innovation as a valuable but risky good, producing health benefits but increasing financial risk for consumers and healthcare payers. This perspective overlooks how innovation can lower physical risks borne by healthy patients facing the prospect of future disease. We present an alternative framework that accounts for all these sources of value and links them to the value of healthcare insurance. We show that any innovation worth buying reduces overall risk and generates positive insurance value on its own. We conduct a stylized numerical exercise to assess the potential empirical significance of our insights. Our calculations suggest that conventional methods meaningfully understate the value of historical health gains and disproportionately undervalue treatments for the most severe illnesses, where physical risk to consumers is the costliest. These calculations also suggest that the value of physical insurance from new technologies may exceed the financial spending risk that they pose.},
   keywords = {Cost-benefit analysis
Health insurance
Medical innovation
Value of health},
   ISSN = {0047-2727},
   DOI = {https://doi.org/10.1016/j.jpubeco.2016.11.012},
   url = {http://www.sciencedirect.com/science/article/pii/S0047272716301864},
   year = {2017},
   type = {Journal Article}
}

@article{ChandraMalani2017civilRightsHealth,
   author = {Chandra, Amitabh and Frakes, Michael and Malani, Anup},
   title = {Challenges To Reducing Discrimination And Health Inequity Through Existing Civil Rights Laws},
   journal = {Health Affairs},
   volume = {36},
   number = {6},
   pages = {1041-1047},
   note = {doi: 10.1377/hlthaff.2016.1091},
   abstract = {More than fifty years after the passage of the Civil Rights Act of 1964, health care for racial and ethnic minorities remains in many ways separate and unequal in the United States. Moreover, efforts to improve minority health care face challenges that differ from those confronted during de jure segregation. We review these challenges and examine whether stronger enforcement of existing civil rights legislation could help overcome them. We conclude that stronger enforcement of existing laws?for example, through executive orders to strengthen enforcement of the laws and congressional action to allow private individuals to bring lawsuits against providers who might have engaged in discrimination?would improve minority health care, but this approach is limited in what it can achieve. Complementary approaches outside the legal arena, such as quality improvement efforts and direct transfers of money to minority-serving providers?those seeing a disproportionate number of minority patients relative to their share of the population?might prove to be more effective.},
   ISSN = {0278-2715},
   DOI = {10.1377/hlthaff.2016.1091},
   url = {https://doi.org/10.1377/hlthaff.2016.1091},
   year = {2017},
   type = {Journal Article}
}

@article{ChandraMalani2017civilRightsHealthReply,
   author = {Chandra, Amitabh and Frakes, Michael and Malani, Anup},
   title = {Civil Rights Laws: The Authors Reply},
   journal = {Health Affairs},
   volume = {36},
   number = {12},
   pages = {2212-2212},
   note = {doi: 10.1377/hlthaff.2017.1407},
   ISSN = {0278-2715},
   DOI = {10.1377/hlthaff.2017.1407},
   url = {https://doi.org/10.1377/hlthaff.2017.1407},
   year = {2017},
   type = {Journal Article}
}


%%%%%%%%%%
%% 2015 %%
%%%%%%%%%%

@article{BennettMalani2015learningSARSTaiwan,
   author = {Bennett, Daniel and Chiang, Chun-Fang and Malani, Anup},
   title = {Learning during a crisis: The SARS epidemic in Taiwan},
   journal = {Journal of Development Economics},
   volume = {112},
   pages = {1-18},
   abstract = {SARS struck Taiwan in 2003, causing a national crisis. Many people feared that SARS would spread through the health care system, and outpatient visits fell by more than 30\% in the course of a few weeks. We examine how both public information and the behavior and opinions of peers contributed to this reaction. We identify a peer effect through a difference-in-difference comparison of longtime residents and recent arrivals, who are less socially connected. Although several forms of social interaction may contribute to this pattern, social learning is a plausible explanation for our finding. We find that people respond to both public information and to their peers. In a dynamic simulation based on the regressions, social interactions substantially magnify the response to SARS.},
   keywords = {SARS
Social learning
Peer effects
Prevalence response
Economic epidemiology
Crisis},
   ISSN = {0304-3878},
   DOI = {https://doi.org/10.1016/j.jdeveco.2014.09.006},
   url = {https://www.sciencedirect.com/science/article/pii/S0304387814001059},
   year = {2015},
   type = {Journal Article}
}

@article{Malani2015anticipationEffects,
   author = {Malani, Anup and Reif, Julian},
   title = {Interpreting pre-trends as anticipation: Impact on estimated treatment effects from tort reform},
   journal = {Journal of Public Economics},
   volume = {124},
   pages = {1-17},
   abstract = {While conducting empirical work, researchers sometimes observe changes in outcomes before adoption of a new policy. The conventional diagnosis is that treatment is endogenous. This observation is also consistent, however, with anticipation effects that arise naturally out of many theoretical models. This paper illustrates that distinguishing endogeneity from anticipation matters greatly when estimating treatment effects. It provides a framework for comparing different methods for estimating anticipation effects and proposes a new set of instrumental variables to address the problem that subjects' expectations are unobservable. Finally, this paper examines a specific set of tort reforms that was not targeted at physicians but was likely anticipated by them. Interpreting pre-trends as evidence of anticipation increases the estimated effect of these reforms by a factor of two compared to a model that ignores anticipation.},
   keywords = {Anticipation
Medical malpractice
Endogeneity
Tort reform},
   ISSN = {0047-2727},
   DOI = {http://doi.org/10.1016/j.jpubeco.2015.01.001},
   url = {http://www.sciencedirect.com/science/article/pii/S0047272715000122},
   year = {2015},
   type = {Journal Article}
}

@article{NandiMalani2015beterEvidenceRSBY,
   author = {Nandi, Arindam and Holtzman, E. Phoebe and Malani, Anup and Laxminarayan, Ramanan},
   title = {The need for better evidence to evaluate the health & economic benefits of India's Rashtriya Swasthya Bima Yojana},
   journal = {Indian Journal of Medical Research},
   volume = {142},
   number = {4},
   pages = {383-390},
   abstract = {In this review the existing evidence on the impact of Rashtriya Swasthya Bima Yojana (RSBY) is discussed in the context of international literature available on health insurance. We describe potential pathways through which health insurance can affect health and economic outcomes, discuss evidence from other developing countries, and identify potential biases and inconsistencies in existing studies on RSBY impact. Given the relatively recent introduction of RSBY, lack of quality, verifiable data on utilization patterns, and the absence of reliable evaluation studies, there is a need to exercise caution while assessing the merits of the programme. Considering the enormous potential and cost of the programme, we emphasize the need for a rigorous impact evaluation of RSBY. It will not only help capture the real impact of the scheme, but may also be able to estimate the extent of systemic inefficiencies at the level of the consumer.},
   keywords = {Financial risk protection
impact evaluation
RSBY
social health insurance},
   ISSN = {0971-5916 (Print)},
   url = {https://journals.lww.com/ijmr/fulltext/2015/42040/the_need_for_better_evidence_to_evaluate_the.6.aspx},
   year = {2015},
   type = {Journal Article}
}


% NON-JOURNAL

@inbook{LakdawallaMalani2015complexHCReformInnovation,
   author = {Lakdawalla, Darius and Malani, Anup and Reif, Julian},
   title = {The Complex Relationship Between Health Care Reform and Innovation},
   booktitle = {The Future of Healthcare Reform in the United States},
   editor = {Malani, Anup and Schill, Michael H.},
   publisher = {University of Chicago Press},
   pages = {289},
   ISBN = {022625495X},
   year = {2015},
   type = {Book Section}
}

@book{RN6554,
   author = {Malani, A. and Schill, M.H.},
   title = {The Future of Healthcare Reform in the United States},
   publisher = {University of Chicago Press},
   ISBN = {9780226255002},
   url = {https://books.google.com/books?id=SJV1CgAAQBAJ},
   year = {2015},
   type = {Edited Book}
}


% WORKING PAPER

@article{BakerMalani2015judgeCareAboutLaw,
   author = {Baker, Scott and Malani, Anup},
   title = {Do judges actually care about the law? Evidence from circuit split data},
   journal = {Conference on Empirical Legal Studies},
   year = {2015},
   type = {Journal Article}
}

%%%%%%%%%%
%% 2014 %%
%%%%%%%%%%

@article{HoldenMalani2014renogiatationByContract,
   author = {Holden, Richard and Malani, Anup},
   title = {Renegotiation Design by Contract},
   journal = {The University of Chicago Law Review},
   volume = {81},
   number = {1},
   pages = {151-178},
   ISSN = {00419494},
   url = {http://www.jstor.org/stable/23646373},
   year = {2014},
   type = {Journal Article}
}

@article{BakerMalani2014enforcerDilemmaCourts,
   author = {Baker, Scott and Malani, Anup},
   title = {Trial Court Budgets, the Enforcer's Dilemma, and the Rule of Law},
   journal = {University of Illinois Law Review},
   volume = {2014},
   number = {5},
   pages = {1573-1602},
   url = {https://heinonline.org/HOL/P?h=hein.journals/unilllr2014&i=1593
https://heinonline.org/HOL/PrintRequest?handle=hein.journals/unilllr2014&collection=journals&div=48&id=1593&print=section&sction=48},
   year = {2014},
   type = {Journal Article}
}

@article{LaxminarayanMalani2014incentiveReportOutbreaksMeningitis,
   author = {Laxminarayan, Ramanan and Reif, Julian and Malani, Anup},
   title = {Incentives for Reporting Disease Outbreaks},
   journal = {PLOS ONE},
   volume = {9},
   number = {3},
   pages = {e90290},
   abstract = {Background Countries face conflicting incentives to report infectious disease outbreaks. Reports of outbreaks can prompt other countries to impose trade and travel restrictions, which has the potential to discourage reporting. However, reports can also bring medical assistance to contain the outbreak, including access to vaccines. Methods We compiled data on reports of meningococcal meningitis to the World Health Organization (WHO) from 54 African countries between 1966 and 2002, a period is marked by two events: first, a large outbreak reported from many countries in 1987 associated with the Hajj that resulted in more stringent requirements for meningitis vaccination among pilgrims; and second, another large outbreak in Sub-Saharan Africa in 1996 that led to a new international mechanism to supply vaccines to countries reporting a meningitis outbreak. We used fixed-effects regression modeling to statistically estimate the effect of external forcing events on the number of countries reporting cases of meningitis to WHO. Findings We find that the Hajj vaccination requirements started in 1988 were associated with reduced reporting, especially among countries with relatively fewer cases reported between 1966 and 1979. After the vaccine provision mechanism was in place in 1996, reporting among countries that had previously not reported meningitis outbreaks increased. Interpretation These results indicate that countries may respond to changing incentives to report outbreaks when they can do so. In the long term, these incentives are likely to be more important than surveillance assistance in prompt reporting of outbreaks.},
   DOI = {10.1371/journal.pone.0090290},
   url = {https://doi.org/10.1371/journal.pone.0090290},
   year = {2014},
   type = {Journal Article}
}

%%%%%%%%%%
%% 2013 %%
%%%%%%%%%%

@article{FarnsworthMalani2013policyPreferenceInterpretation,
   author = {Farnsworth, Ward and Guzior, Dustin and Malani, Anup},
   title = {Policy Preferences and Legal Interpretation},
   journal = {Journal of Law and Courts},
   volume = {1},
   number = {1},
   pages = {115-138},
   abstract = {This article presents an empirical study of statutory interpretation. Respondents were asked to read statutes and answer questions about how they should be applied to simple cases. The results suggest, first, that it is difficult to separate judgments about the linguistic meaning of a statute from policy preferences about it. Different ways of framing the interpretive question have consequences, however; asking how an ordinary reader would interpret a text helps produce answers that are distinct from a respondent’s own preferences about it. The article considers why this might be so and discusses implications for the interpretation of statutes by courts.},
   ISSN = {2164-6570},
   DOI = {10.1086/668603},
   url = {https://www.cambridge.org/core/product/F8DBE7EB48916F39F720942A96B76008},
   year = {2013},
   type = {Journal Article}
}

@article{BoniMalani2013econEpidemiologyFlu,
   author = {Boni, Maciej F. and Galvani, Alison P. and Wickelgren, Abraham L. and Malani, Anup},
   title = {Economic epidemiology of avian influenza on smallholder poultry farms},
   journal = {Theoretical Population Biology},
   volume = {90},
   pages = {135-144},
   abstract = {Highly pathogenic avian influenza (HPAI) is often controlled through culling of poultry. Compensating farmers for culled chickens or ducks facilitates effective culling and control of HPAI. However, ensuing price shifts can create incentives that alter the disease dynamics of HPAI. Farmers control certain aspects of the dynamics by setting a farm size, implementing infection control measures, and determining the age at which poultry are sent to market. Their decisions can be influenced by the market price of poultry which can, in turn, be set by policy makers during an HPAI outbreak. Here, we integrate these economic considerations into an epidemiological model in which epidemiological parameters are determined by an outside agent (the farmer) to maximize profit from poultry sales. Our model exhibits a diversity of behaviors which are sensitive to (i) the ability to identify infected poultry, (ii) the average price of infected poultry, (iii) the basic reproductive number of avian influenza, (iv) the effect of culling on the market price of poultry, (v) the effect of market price on farm size, and (vi) the effect of poultry density on disease transmission. We find that under certain market and epidemiological conditions, culling can increase farm size and the total number of HPAI infections. Our model helps to inform the optimization of public health outcomes that best weigh the balance between public health risk and beneficial economic outcomes for farmers.},
   keywords = {Avian influenza
H5N1
Economic epidemiology
Culling},
   ISSN = {0040-5809},
   DOI = {https://doi.org/10.1016/j.tpb.2013.10.001},
   url = {https://www.sciencedirect.com/science/article/pii/S0040580913001044},
   year = {2013},
   type = {Journal Article}
}

@article{KamenicaMalani2013placeboEffectAds,
   author = {Kamenica, Emir and Naclerio, Robert and Malani, Anup},
   title = {Advertisements impact the physiological efficacy of a branded drug},
   journal = {Proceedings of the National Academy of Sciences},
   volume = {110},
   number = {32},
   pages = {12931},
   abstract = {We conducted randomized clinical trials to examine the impact of direct-to-consumer advertisements on the efficacy of a branded drug. We compared the objectively measured, physiological effect of Claritin (Merck &amp; Co.), a leading antihistamine medication, across subjects randomized to watch a movie spliced with advertisements for Claritin or advertisements for Zyrtec (McNeil), a competitor antihistamine. Among subjects who test negative for common allergies, exposure to Claritin advertisements rather than Zyrtec advertisements increases the efficacy of Claritin. We conclude that branded drugs can interact with exposure to television advertisements.},
   DOI = {10.1073/pnas.1012818110},
   url = {http://www.pnas.org/content/110/32/12931.abstract},
   year = {2013},
   type = {Journal Article}
}

% NON-JOURNAL

@article{BhattacharyaMalani2013bestBothWorldsHealthcareReform,
   author = {Bhattacharya, Jayanta and Chandra, Amitabh and Chernew, Michael and Goldman, Dana and Jena, Anupam and Lakdawalla, Darius and Malani, Anup and Philipson, Tomas},
   title = {Best of Both Worlds: Uniting Universal Coverage and Personal Choice in Health Care},
   journal = {American Enterprise Institute},
   year = {2013},
   type = {Journal Article}
}


%%%%%%%%%%
%% 2012 %%
%%%%%%%%%%

@article{Malani2012patentDamages,
   author = {Malani, Anup and Masur, Jonathan S.},
   title = {Raising the Stakes in Patent Cases},
   journal = {Georgetown Law Journal},
   volume = {101},
   number = {3},
   pages = {637-688},
   url = {https://heinonline.org/HOL/LandingPage?handle=hein.journals/glj101&div=21&id=&page=},
   year = {2012},
   type = {Journal Article}
}

@article{Malani2012hetergeneousTreatmentEffects,
   author = {Malani, Anup and Bembom, Oliver and Van der Laan, Mark},
   title = {Accounting for Heterogeneous Treatment Effects in the FDA Approval Process},
   journal = {Food and Drug Law Journal},
   volume = {67},
   number = {1},
   pages = {23-50},
   url = {https://heinonline.org/HOL/P?h=hein.journals/foodlj67&i=31},
   year = {2012},
   type = {Journal Article}
}

@article{YangMalani2012decisionTriangleBehavioralEcon,
   author = {Yang, Haiyang and Carmon, Ziv and Kahn, Barbara and Malani, Anup and Schwartz, Janet and Volpp, Kevin and Wansink, Brian},
   title = {The Hot–Cold Decision Triangle: A framework for healthier choices},
   journal = {Marketing Letters},
   volume = {23},
   number = {2},
   pages = {457-472},
   abstract = {People often behave in ways that are clearly detrimental to their health. We review representative research on unhealthy behaviors within a parsimonious framework, the Hot-Cold Decision Triangle. Through this framework, we describe how when people embrace colder state reasoning—instead of risking the pitfalls of heuristics and visceral reactions—they are more likely to behave healthily. We also illustrate how some heuristics and visceral urges can be leveraged to encourage healthier choices. We conclude by discussing unexplored research directions, as well as substantive implications for individuals, marketers, and policymakers.},
   ISSN = {1573-059X},
   DOI = {10.1007/s11002-012-9179-0},
   url = {https://doi.org/10.1007/s11002-012-9179-0},
   year = {2012},
   type = {Journal Article}
}


% BOOK

@inbook{Malani2012regMedicalProductsHandbook,
   author = {Malani, Anup and Philipson, Tomas J.},
   title = {The Regulation of Medical Products},
   booktitle = {The Oxford Handbook of the Economics of the Biopharmaceutical Industry},
   editor = {Danzon, Patricia M. and Nicholson, Sean},
   publisher = {Oxford University Press},
   pages = {100-142},
   abstract = {This article synthesizes and extends, in a nontechnical manner, recent research on the Food and Drug Administration (FDA). The aim is to shed light on whether the policies of the agency itself are safe and effective when measured in terms of economic efficiency. The first section provides an overview of the role of the FDA in regulating pharmaceutical drugs and medical devices. The second section surveys the existing efficiency rationales for government regulation of the information about and the quality of medical products, and then canvasses the literature for empirical studies on the effects of FDA regulation on innovation and costs. The final section examines the growing role of tort law—specifically, products liability litigation—in supplementing FDA regulation of drug quality.},
   ISBN = {9780199742998},
   DOI = {10.1093/oxfordhb/9780199742998.013.0005},
   url = {https://doi.org/10.1093/oxfordhb/9780199742998.013.0005},
   year = {2012},
   type = {Book Section}
}

@inbook{Malani2012commentCohenAntiMalarialDiagnostics,
   author = {Malani, Anup},
   title = {Commentary on: Jessica L. Cohen and William T. Dickens, Adoption of Over-the-Counter Malaria Diagnostics in Africa: The Role of Subsidies, Beliefs, Externalities, and Competition},
   booktitle = {The Value of Information: Methodological Frontiers and New Applications in Environment and Health},
   editor = {Laxminarayan, Ramanan and Macauley, Molly K.},
   publisher = {Springer Netherlands},
   address = {Dordrecht},
   pages = {188-191},
   abstract = {Plans for the wide-scale distribution and subsidy of artemisinin combination therapies (ACTs), an antimalarial treatment, pose two problems for public health planning. First, many people seeking malaria treatment do not have the disease. If ACT subsidies could be targeted toward those with malaria, the cost of subsidies could fall. Second, the inappropriate use of antimalarial drugs may contribute to the emergence of drug-resistant parasites. Rapid diagnostic tests (RDTs) for malaria could help with both problems, but drug shop owners may have few financial incentives to sell them, given profits from overtreatment for malaria. A model of the provision of RDTs by profit-maximizing drug shops shows that if all parties know the probability of having malaria and if there are no subsidies for drugs and no external costs to inappropriate treatment, both monopolistic and competitive drug shop owners will provide RDTs under the same circumstances that a social welfare maximizing planner would. However, since drugs will be subsidized, customers overestimate their likelihood of having malaria, and since there are external costs to the misuse of antimalarials, profit-maximizing drug shops will likely underprovide RDTs. We show that a subsidy for RDTs can increase provision and, under adequate competition, induce everyone to use RDTs optimally. The results also highlight the importance of educating customers about the true prevalence of malaria and promoting competition among drug providers.},
   ISBN = {978-94-007-4839-2},
   DOI = {10.1007/978-94-007-4839-2_7},
   url = {https://doi.org/10.1007/978-94-007-4839-2_7},
   year = {2012},
   type = {Book Section}
}


%%%%%%%%%%
%% 2011 %%
%%%%%%%%%%

@article{Malani2011incentiveReportingDisease,
   author = {Malani, Anup and Laxminarayan, Ramanan},
   title = {Incentives for Reporting Infectious Disease Outbreaks},
   journal = {Journal of Human Resources},
   volume = {46},
   number = {1},
   pages = {176},
   abstract = {The global spread of diseases such as swine flu and SARS highlights the difficult decision governments face when presented with evidence of a local outbreak. Reporting the outbreak may bring medical assistance but is also likely to trigger trade sanctions by countries hoping to contain the disease. Suppressing the information may avoid trade sanctions, but increases the likelihood of widespread epidemics. In this paper, we model the government’s decision as a signaling game in which a country has private but imperfect evidence of an outbreak. First, we find that not all sanctions discourage reporting. Sanctions based on fears of an undetected outbreak (false negatives) encourage disclosure by reducing the relative cost of sanctions that follow a reported outbreak. Second, improving the quality of detection technology may not promote the disclosure of an outbreak because the forgone trade from reporting truthfully is that much greater. Third, informal surveillance is an important channel for publicizing outbreaks and functions as an exogenous yet imperfect signal that is less likely to discourage disclosure. In sum, obtaining accurate information about potential epidemics is as much about reporting incentives as it is about detection technology.},
   DOI = {10.3368/jhr.46.1.176},
   url = {http://jhr.uwpress.org/content/46/1/176.abstract},
   year = {2011},
   type = {Journal Article}
}

% NON-JOURNAL

@misc{LaxminarayanMalani2011handbookEconomicsInfectiousDisease,
   author = {Laxminarayan, Ramanan and Malani, Anup},
   title = {189 Economics of Infectious Diseases},
   publisher = {Oxford University Press},
   pages = {0},
   abstract = {Infectious diseases remain a central preoccupation in many countries. This article sets out the economic issues that arise in this highly complex domain. It reviews four main strands of literature on the economics of infectious diseases. It discusses the economic impact of infectious diseases on labor productivity and investment decisions. It focuses on the interplay between disease prevention and treatment and individual risk-taking behavior. It discusses vaccination as an important tool in the prevention of infectious diseases and presents a classic public goods problem. Disease reporting and eradication efforts are also global public goods. A fourth strand of literature is on the optimal design and allocation of resources for prevention and treatment programs. These programs are based on epidemiological models of disease spread that present significant mathematical challenges.},
   ISBN = {9780199238828},
   DOI = {10.1093/oxfordhb/9780199238828.013.0009},
   url = {https://doi.org/10.1093/oxfordhb/9780199238828.013.0009},
   year = {2011},
   type = {Electronic Book Section}
}



%%%%%%%%%%
%% 2010 %%
%%%%%%%%%%


@article{FarnsworthMalani2010ambiguity,
   author = {Farnsworth, Ward and Guzior, Dustin F. and Malani, Anup},
   title = {Ambiguity about Ambiguity: An Empirical Inquiry into Legal Interpretation},
   journal = {Journal of Legal Analysis},
   volume = {2},
   number = {1},
   pages = {257-300},
   abstract = {Most scholarship on statutory interpretation discusses what courts should do with ambiguous statutes. This paper investigates the crucial and analytically prior question of what ambiguity in law is. Does a claim that a text is ambiguous mean the judge is uncertain about its meaning? Or is it a claim that ordinary readers of English, as a group, would disagree about what the text means? This distinction is of considerable theoretical interest. It also turns out to be highly consequential as a practical matter.To demonstrate, we developed a survey instrument for exploring determinations of ambiguity and administered it to nearly 1,000 law students. We find that asking respondents whether a statute is “ambiguous” in their own minds produces answers that are strongly biased by their policy preferences. But asking respondents whether the text would likely be read the same way by ordinary readers of English does not produce answers biased in this way. This discrepancy leads to important questions about which of those two ways of thinking about ambiguity is more legally relevant. It also has potential implications for how cases are decided and for how law is taught.},
   ISSN = {2161-7201},
   DOI = {10.1093/jla/2.1.257},
   url = {https://doi.org/10.1093/jla/2.1.257},
   year = {2010},
   type = {Journal Article}
}

@article{EberMalani2010hospitalAcquiredInfections,
   author = {Eber, Michael R. and Laxminarayan, Ramanan and Perencevich, Eli N. and Malani, Anup},
   title = {Clinical and Economic Outcomes Attributable to Health Care–Associated Sepsis and Pneumonia},
   journal = {Archives of Internal Medicine},
   volume = {170},
   number = {4},
   pages = {347-353},
   abstract = {Health care–associated infections affect 1.7 million hospitalizations each year, but the clinical and economic costs attributable to these infections are poorly understood. Reliable estimates of these costs are needed to efficiently target limited resources for the greatest public health benefit.Hospital discharge records from the Nationwide Inpatient Sample database were used to identify sepsis and pneumonia cases among 69 million discharges from hospitals in 40 US states between 1998 and 2006. Community-acquired infections were excluded using criteria adapted from previous studies. Because these criteria may not exclude all community-acquired infections, outcomes were examined separately for cases associated with invasive procedures, which were unlikely to result from preexisting infections. Attributable hospital length of stay, hospital costs, and crude in-hospital mortality were estimated from discharge records using a multivariate matching analysis and a supplementary regression analysis. These models controlled for patient diagnoses, procedures, comorbidities, demographics, and length of stay before infection.In cases associated with invasive surgery, attributable mean length of stay was 10.9 days, costs were $32,900, and mortality was 19.5\% for sepsis; corresponding values for pneumonia were 14.0 days, $46,400, and 11.4\%, respectively (P$<$.001). In cases not associated with invasive surgery, attributable mean length of stay, costs, and mortality were estimated to be 1.9 to 6.0 days, $5800 to $12,700, and 11.7\% to 16.0\% for sepsis and 3.7 to 9.7 days, $11,100 to $22,300, and 4.6\% to 10.3\% for pneumonia (P$<$.001).Health care–associated sepsis and pneumonia impose substantial clinical and economic costs.Arch Intern Med. 2010;170(4):347-353-->},
   ISSN = {0003-9926},
   DOI = {10.1001/archinternmed.2009.509},
   url = {https://doi.org/10.1001/archinternmed.2009.509},
   year = {2010},
   type = {Journal Article}
}


